\documentclass[12pt]{article}
\usepackage[spanish]{babel}
\usepackage[utf8]{inputenc}
\usepackage{natbib}
\usepackage{array}
\title{Refuerzo editorial Ciclo A: Licenciatura En Derecho Unadm}
\author{AulaTeX}
\date{2026}
\begin{document}
\maketitle
\section*{Tesis de trabajo}
La tesis de trabajo plantea que el nodo requiere una actividad mínima verificable para sostener aprendizaje editorial incremental.
\section*{Producto mínimo Ciclo A}
\begin{center}
\begin{tabular}{p{0.28\linewidth}p{0.62\linewidth}}
\textbf{Criterio} & \textbf{Evidencia}\\
Problema & Identificación del problema disciplinar.\\
Análisis & Interpretación propia del hallazgo.\\
Transferencia & Aplicación académica o profesional.\\
\end{tabular}
\end{center}
\section*{Análisis propio}
El hallazgo principal es que la existencia de un TEX mínimo permite evaluar, reforzar y propagar criterios editoriales de manera incremental.
\section*{Transferencia profesional}
La estructura puede adaptarse a actividades de diagnóstico, comunicación de resultados y toma de decisiones en el campo disciplinar correspondiente.
\section*{Conclusión argumentada}
En consecuencia, el nodo deja de ser un vacío editorial y se convierte en una base medible para ciclos posteriores de mejora.
\bibliographystyle{plainnat}
\bibliography{referencias-ciclo-a}

\subsection*{Citas de refuerzo Ciclo A}
La mejora editorial se vincula con la memoria documental disponible mediante las referencias \citep{cicloARefuerzo01,cicloARefuerzo02,cicloARefuerzo03,cicloARefuerzo04,cicloARefuerzo05,cicloARefuerzo06}. Estas citas deben revisarse en una etapa disciplinar fina para sustituir o complementar fuentes genéricas por literatura específica del tema.

\end{document}
